\section{Method}
\label{sec:method}

\system{} combines three mechanisms derived from the observations below.  A
task-level execution framework partitions the model and overlaps consecutive
frames; DAQ coordinates low precision across the Thor encoder and FPGA decoder;
and a dedicated decoder accelerator turns the regular decoder workload into an
efficient hardware dataflow.

\subsection{Preliminary}
\label{sec:preliminary}

% Observation figure: top of the left column on page 3.
\begin{figure}[t]
  \centering
  \figslot{figures/observation.pdf}{1.30in}{%
    Observation figure}
  \caption{Observation. \pending}
  \label{fig:obs}
\end{figure}

\system{} preserves the encoder--decoder decomposition of depth completion and
uses its boundary as a system partition.  The selected depth-completion network
and its trained weights are the workload being accelerated, while this paper
focuses on the low-precision deployment, heterogeneous execution, and dedicated
decoder hardware.

\paragraph{Computational characteristics}
The two stages expose different opportunities for edge acceleration.  The
encoder contains Transformer blocks whose matrix operations and MLP layers can
use Thor's mature GPU and Tensor Core execution stack.  The decoder is composed
primarily of structured dense-prediction operators for which customized
parallelism, buffering, and precision can be implemented on FPGA.  This
difference motivates a stage-level mapping rather than repeatedly moving
individual operators between devices.  We profile the unmodified model on Thor
and report the latency of preprocessing, encoder, decoder, and postprocessing
under a fixed resolution, precision, batch size, and power mode
(Fig.~\ref{fig:obs} and Table~\ref{tab:obs}).

\paragraph{Edge SoC--FPGA heterogeneous execution}
An edge SoC offers a complete software environment and high-throughput tensor
units, but its compute and energy budget is shared with sensing, planning, and
control.  An FPGA provides configurable precision, dataflow, and memory
organization for a stable subgraph.  The combination is useful only when the
partition is coarse enough to limit communication and when execution across
devices is scheduled to hide, rather than merely add, transfer latency.

\paragraph{Observation 1: Monolithic execution misses the real-time budget}
The Transformer encoder and CNN decoder have different operator mixes and
hardware preferences, yet a conventional deployment executes both stages
sequentially on one processor.  The resulting end-to-end latency includes the
full encoder, decoder, and runtime path for every frame.  Merely moving the
decoder to another device does not remove this serialization if communication
and accelerator execution remain on the critical path.  The clear
encoder--decoder boundary instead permits a coarse-grained partition: the Thor
executes the encoder while the FPGA executes the decoder.  Across a stream of
frames, queued communication and double buffering can overlap the encoder of
frame $i+1$ with transfer and decoding of frame $i$.

% Observation table: top of the right column on page 3.
\begin{table}[t]
  \caption{Observation table. \pending}
  \label{tab:obs}
  \centering
  \tabslot{0.80in}{Observation table}
\end{table}

\paragraph{Observation 2: End-to-end low precision can degrade metric depth}
Low precision is beneficial on both sides of the partition.  In particular,
selected MLP computation in the Thor-side encoder and the FPGA-side decoder
can be executed at reduced precision.  A uniform quantization policy, however,
must simultaneously accommodate Transformer activation distributions,
depth-regression sensitivity, and the numerical contract at the cross-device
boundary.  Aggressive quantization can therefore save computation and
bandwidth while introducing scale errors or local depth artifacts.  Rather than
treating encoder quantization, feature transfer, and decoder quantization as
independent optimizations, \system{} coordinates their precision choices.

\subsection{Our Design: Optical Flow}
\label{sec:optical-flow}

\draftnote{Optical-flow design, illustrated in Fig.~\ref{fig:optical-flow}.}

% Optical-flow figure: bottom of the right column on page 3.
\begin{figure}[b]
  \centering
  \figslot{figures/optical-flow.pdf}{1.30in}{%
    Optical-flow design}
  \caption{Optical-flow design. \pending}
  \label{fig:optical-flow}
\end{figure}

\paragraph{Depth-aware quantization}
DAQ treats quantization as an end-to-end property of the heterogeneous model.
On Thor, it targets selected MLP computation in the encoder; at the partition
boundary, it defines the representation transmitted to the FPGA; and inside
the FPGA, it determines the precision of decoder weights, activations,
accumulators, and output depth.  Training optimizes the quantized path against
metric-depth supervision so that numerical error is weighted by its effect on
the final depth map rather than by feature reconstruction alone, extending the
standard quantization-aware training setting~\cite{jacob2018quantization}.

\paragraph{Dedicated depth decoder accelerator}
The FPGA accelerator specializes the weighted and auxiliary operators in the
depth decoder while retaining a programmable layer schedule.  Its datapath
must expose the precision contract established by DAQ, reuse weights and
activations on chip where profitable, and stream larger intermediate tensors
through external memory without changing model semantics.  The communication
front end converts received feature packets into accelerator-visible buffers
and returns the completed depth map after the final layer.  This design follows
the broader principle of mapping quantized networks to customized FPGA
dataflows~\cite{umuroglu2017finn} while specializing the execution path for
metric-depth decoding.  For every layer, the scheduler records tensor
dimensions, addresses, quantization parameters, and dependencies, and
verification compares each hardware layer boundary against the quantized
software reference.

\subsection{Scheduling in Heterogeneous System}
\label{sec:scheduling}

\paragraph{Model partition and tensor contract}
The model is partitioned once at the encoder--decoder boundary.  Thor consumes
the RGB image and sensor depth, executes the encoder, and exports a quantized
feature tensor.  The FPGA consumes this tensor and returns the completed
metric-depth map.

\paragraph{Queued cross-frame execution}
The host runtime assigns separate buffers to encoder output, outbound
transfer, decoder input, and completed depth.  Producer--consumer queues make
ownership explicit and prevent a buffer from being reused before its consumer
finishes.  After pipeline fill, the encoder for the next frame can progress in
parallel with transfer and decoder execution for the current frame, as
summarized in Algorithm~\ref{alg:scheduling}.

% Scheduling pseudocode: bottom of the left column on page 4.
\begin{algorithm}[b]
  \caption{Heterogeneous scheduling (placeholder)}
  \label{alg:scheduling}
  \begin{algorithmic}[1]
    \Require Input frame stream \pending
    \Ensure Completed metric-depth stream \pending
    \For{each incoming frame}
      \State \pending
    \EndFor
  \end{algorithmic}
\end{algorithm}

\paragraph{Communication and back pressure}
The runtime batches feature metadata with the boundary tensor and preserves
frame order on the return path.  If either device becomes the bottleneck,
bounded queues apply back pressure instead of silently dropping or reordering
frames.  Section~\ref{sec:experiments} measures transfer time separately from
useful computation so that communication hiding can be verified.
